\chapter{Test}

\section{Strategia di verifica}

La verifica del progetto si articola su tre livelli: test unitari della logica applicativa, test di integrazione sugli endpoint e test end-to-end sui flussi principali. Questa combinazione e' coerente con un'applicazione distribuita ma con dominio applicativo ben separato.

\section{Test backend}

Nel repository corrente sono presenti suite di test per ogni microservizio backend, con focus su use case e router HTTP. Le principali aree coperte sono:

\begin{itemize}
\item autenticazione e sessioni;
\item catalogo e operazioni admin;
\item CAD e logica di derivazione BOM;
\item carrello e checkout;
\item notifiche e subscription;
\item chatbot e messaggistica contestuale;
\item ordini e aggiornamento stato;
\item reporting e trend ordini.
\end{itemize}

Nel CAD, ad esempio, i test verificano la geometria delle spine, la derivazione della BOM, l'accettazione o il rifiuto dei livelli, la validazione del piano colonne e il comportamento dei servizi HTTP. Questo e' il punto piu' delicato del sistema, perche' da qui dipende la correttezza di carrello, notifiche e chat contestuale.

\section{Test frontend}

Il frontend dispone di una struttura di composables e store che rende naturale testare i flussi principali per vista. Anche se la copertura automatica della UI non e' ancora esaustiva, il progetto include verifiche di stato su:

\begin{itemize}
\item auth store;
\item cart badge e refresh contatore;
\item notifiche non lette e toast;
\item composables del CAD per list/detail/update/finalize;
\item integrazione con i service client REST e realtime.
\end{itemize}

Le viste principali sono state inoltre progettate con stati di caricamento, messaggi di errore e attributi accessibili che rendono piu' semplice una futura automatizzazione con test UI dedicati.

\section{Test end-to-end}

Nel repository sono presenti suite e2e per i flussi principali:

\begin{itemize}
\item auth integration;
\item cad integration;
\item cart integration;
\item catalog integration;
\item chatbot integration;
\item notifications integration;
\item order integration;
\item reporting integration.
\end{itemize}

Lo workspace contiene anche un comando aggregato \texttt{npm run test:baseline} che esegue la suite backend e la suite e2e, oltre ai comandi separati per ogni microservizio.

\section{Lista sintetica dei test fatti}

Nel corso della documentazione e dell'allineamento del progetto sono stati considerati come riferimenti i seguenti test realmente presenti nel codice:

\begin{itemize}
\item backend: test di unit\`a e router per ciascun microservizio;
\item e2e: suite per auth, CAD, cart, catalog, chatbot, notifiche, ordini e reporting;
\item frontend: verifica della logica di store, composables e service layer;
\item gateway: test dei routing BFF e health route;
\item CAD: test su \texttt{SpineModel}, \texttt{deriveBom}, comando di configurazione e router.
\end{itemize}

\section{Stato della copertura}

La traccia di copertura presente in \texttt{utilities/traceability-matrix.md} mostra che i requisiti principali sono coperti, mentre restano aree di completamento soprattutto sul fronte UX, smoke test accessibilita' e robustezza di alcuni casi realtime. Il documento di progetto deve quindi riflettere un sistema gia' solido sui flussi core, ma con margini di crescita ancora evidenti sul lato UI e validazione finale.
